\structural{ВСТУП}

Інформаційні системи, що обробляють потік завдань різних типів на спільних обчислювальних
ресурсах, мають вирішувати, у якому порядку обслуговувати завдання, що очікують. Від стратегії
упорядкування черги залежать час відповіді системи, кількість завдань, що накопичуються, і
розподіл очікування між типами завдань. Аналітичні методи теорії масового обслуговування дають
точні результати лише для простих систем, а в системі, де завдання займають різну частку
ресурсу, одне з них монополізує систему, а пріоритет вимагає резервувати ресурс, точний
аналітичний розвʼязок отримати складно. Експерименти з реальною системою дорогі й небезпечні
для її користувачів. Тому актуальним методом дослідження є імітаційне моделювання, яке дає
змогу відтворити функціонування системи з урахуванням усіх правил і оцінити вплив стратегій
до їх впровадження.

Тема курсової роботи -- «\kpiTitle». Обʼєкт дослідження -- інформаційна система, яка обробляє
завдання трьох типів (A, B, C), що конкурують за оперативну памʼять: завдання A і~B можуть
виконуватися одночасно, а завдання~C монополізує систему. Предмет дослідження -- вплив стратегії
упорядкування черги на показники обслуговування.

Мета роботи -- визначити вплив пʼяти стратегій упорядкування черги за різних інтенсивностей
надходження завдань на середній час перебування завдання в системі, середню кількість завдань у
черзі та середній час очікування в черзі, а також запропонувати шляхи поліпшення
функціонування системи.

Для досягнення мети розроблено концептуальну модель системи, її формалізовано стохастичною
мережею Петрі~\cite{stetsenko-kursova, stetsenko-modeling}, у якій обмеження на одночасне
виконання завдань задано кратністю дуг, а стратегії упорядкування -- пріоритетами та
ймовірностями запуску переходів. Алгоритм імітації мережі Петрі реалізовано мовою
Java~\cite{java-jls} на основі ядра бібліотеки PetriObjModelPaint~\cite{petriobjmodelpaint},
розширеного засобами обʼєктно-орієнтованого програмування; збірку виконує Apache
Maven~\cite{maven}, тести -- JUnit~5~\cite{junit}. Коректність моделі перевірено верифікацією:
контролем інваріантів мережі, аналізом реакції відгуків на зміну вхідних змінних і порівнянням з
аналітичними результатами для систем M/M/1 і M/M/2~\cite{kleinrock}. Експериментальне дослідження виконано за повним факторним планом
«стратегія $\times$ інтенсивність надходження» з визначенням перехідного періоду, часу
моделювання й кількості прогонів; результати оброблено дисперсійним і регресійним
аналізом мовою Python з бібліотекою SciPy~\cite{scipy}. Пропозицію щодо поліпшення системи оцінено пробним запуском
моделі.
